\documentclass[11pt,a4paper]{article}
\usepackage[utf8]{inputenc}
\usepackage[T1]{fontenc}
\usepackage{lmodern}
\usepackage{geometry}
\geometry{margin=2cm}
\usepackage{amsmath,amssymb}
\usepackage{booktabs}
\usepackage{tabularx}
\usepackage{xcolor}
\usepackage{enumitem}
\usepackage{hyperref}

\hypersetup{
    colorlinks=true,
    linkcolor=blue!70!black,
    urlcolor=teal!70!black,
    citecolor=purple!70!black
}

\setlist[itemize]{noitemsep, topsep=2pt}
\setlist[enumerate]{noitemsep, topsep=2pt}

\title{\textbf{\Large Reporte Ejecutivo: Estudio de Estrellas RR Lyrae en Rubin DP1}\\
\large Síntesis Estructurada de Ngeow \& Bhardwaj (2026, \textit{The Astronomical Journal}, 172:3)}
\author{\textbf{Proyecto Vera C. Rubin} \quad|\quad Resumen Técnico}
\date{Noviembre 2026}

\begin{document}

\maketitle

\begin{abstract}
\noindent Este documento condensa los puntos metodológicos, observacionales y teóricos más relevantes del artículo \textit{“Exploring RR Lyrae Variable Stars in the Vera C. Rubin Observatory Data Preview 1”} (Ngeow \& Bhardwaj, 2026). El objetivo es disponer de una referencia rápida y organizada de cada sección del artículo, detallando los productos de datos empleados, la estrategia de ajuste por plantillas, los resultados de metalicidad y módulos de distancia, y las implicaciones físicas para el sondeo LSST.
\end{abstract}

\vspace{0.5em}

\section*{Ficha Bibliográfica y Parámetros del Trabajo}
\begin{itemize}
    \item \textbf{Título:} \textit{Exploring RR Lyrae Variable Stars in the Vera C. Rubin Observatory Data Preview 1}
    \item \textbf{Autores:} Chow-Choong Ngeow (National Central University, Taiwán) y Anupam Bhardwaj (IUCAA, India).
    \item \textbf{Publicación:} \textit{The Astronomical Journal}, 172:3 (11pp), Julio 2026. DOI: \texttt{10.3847/1538-3881/ae6807}.
    \item \textbf{Conjunto de datos evaluado:} Data Preview 1 (DP1) del Observatorio Vera C. Rubin (ComCam, 48 noches a finales de 2024, filtros $u, g, r, i, z, y$).
    \item \textbf{Modelos teóricos de comparación:} Modelos de pulsación hidrodinámica no lineal de Marconi et al. (2022) adaptados para los filtros de LSST.
\end{itemize}

\vspace{0.5em}
\hrule
\vspace{0.8em}

\section{Sección 1: Introducción y Motivación}
\begin{itemize}
    \item \textbf{Contexto observacional de DP1:} Rubin DP1 se liberó el 30 de junio de 2025 con datos de comisionamiento de la LSSTComCam (mosaico $3\times3$ CCDs, campo $\sim1.3^\circ \times 1.3^\circ$) cubriendo series temporales en 7 campos preseleccionados.
    \item \textbf{Importancia astrofísica de las RR Lyrae:} Estrellas pulsantes viejas (Población II) de la Rama Horizontal (HB). Son candelas estándar de distancia primarias y trazadoras clave de la cinemática/química del halo de la Vía Láctea y galaxias enanas satélites.
    \item \textbf{Limitaciones de estudios previos:} Trabajos iniciales (Malanchev et al. 2025; Choi et al. 2025) solo recuperaron de forma exploratoria un puñado de variables conocidas en DP1 sin comparaciones teóricas rigurosas.
    \item \textbf{Objetivo central del paper:} Realizar un análisis sistemático y cuantitativo de todas las RR Lyrae conocidas en DP1, contrastando sus curvas de luz reales por primera vez con los modelos teóricos para filtros LSST de Marconi et al. (2022).
\end{itemize}

\section{Sección 2: Muestra de RR Lyrae Conocidas en Campos DP1}
\begin{itemize}
    \item \textbf{Plataforma LSDB y Cruce Catálogo (radio de búsqueda $1''$):}
    \begin{itemize}
        \item Se cruzaron las posiciones de DP1 con la base de datos de estrellas variables VSX (AAVSO) y catálogos de la literatura.
        \item Se evaluaron dos catálogos principales de DP1: \texttt{Object} (imágenes co-sumadas directas) y \texttt{DiaObject} (análisis por diferencia de imágenes, DIA).
    \end{itemize}
    \item \textbf{Comportamiento según la densidad estelar:}
    \begin{itemize}
        \item \textit{Campos poco poblados (EDFS, Rubin95, Rubin38):} Recuperación idéntica en \texttt{Object} y \texttt{DiaObject}.
        \item \textit{Campos abarrotados (47Tuc y Fornax):} \texttt{DiaObject} recupera significativamente más estrellas porque la resta de imágenes elimina el fondo estelar saturado y el efecto de mezcla (\textit{blending}).
        \item \textit{Aclaración sobre 47Tuc:} Las RR Lyrae observadas en esta línea de visión no pertenecen al cúmulo globular 47 Tuc, sino a la \textbf{Pequeña Nube de Magallanes (SMC)} de fondo.
        \item \textit{Campo ECDFS:} Pese a tener la mayor cadencia (21 noches), no posee ninguna RR Lyrae conocida.
    \end{itemize}
    \item \textbf{Incompletitud en Fornax y Muestra DES:} VSX arrojó solo 3 a 10 estrellas en Fornax. Al cruzar con el catálogo de Stringer et al. (2021) del \textit{Dark Energy Survey} (DES), la muestra aumentó a \textbf{393 en \texttt{Object} y 617 en \texttt{DiaObject}}.
    \item \textbf{Periodos y Extinción adoptados:} Se adoptaron los periodos ya conocidos de la literatura (no se ajustaron periodos con DP1). La extinción $E(B-V)$ se tomó de los mapas de Schlegel et al. (1998) y coeficientes $R_\lambda$ de Schlafly \& Finkbeiner (2011) ($R_V = 3.1$).
\end{itemize}

\begin{table}[h]
\centering
\small
\caption{Resumen de la muestra de RR Lyrae identificadas en DP1 (Tabla 1 del artículo)}
\begin{tabularx}{\textwidth}{lcccccX}
\toprule
\textbf{Campo} & \textbf{Épocas} & $N_{OBJ}$ & $N_{DIA}$ & $N_{Final}$ (OBJ / DIA) & \textbf{Tipo de Muestra / Población} \\
\midrule
47Tuc    & 4  & 25  & 62  & 1 / 36   & VSX (Población de la SMC de fondo) \\
ECDFS    & 21 & 0   & 0   & 0 / 0    & VSX (Sin RR Lyrae conocidas) \\
EDFS     & 9  & 1   & 1   & 1 / 1    & VSX (1 RRab de prueba Euclid/Rubin) \\
Fornax   & 2  & 3   & 10  & 0 / 7    & VSX (Incompleto) \\
Rubin95  & 10 & 2   & 2   & 2 / 2    & VSX (1 RRab y 1 RRc bien muestreadas) \\
Rubin38  & 5  & 4   & 4   & 1 / 1    & VSX (1 RRab con calidad suficiente) \\
Seagull  & 4  & 1   & 1   & 0 / 0    & VSX (Descartado por alta dispersión) \\
\textbf{Fornax (DES)} & 2 & \textbf{393} & \textbf{617} & \textbf{146 / 517} & Stringer et al. (2021) / DES \\
\bottomrule
\end{tabularx}
\end{table}

\section{Sección 3: Curvas de Luz y Ajuste de Plantillas}
\begin{itemize}
    \item \textbf{Selección de fotometría forzada:} Se seleccionó \texttt{diaObjectForcedSource} como producto primario de DIA (siguiendo a Choi et al. 2025), ya que mide flujo forzado en todas las visitas (a diferencia de \texttt{diaSource}, que pierde épocas cuando la estrella está cerca del mínimo de brillo).
    \item \textbf{Filtrado de calidad estricto:}
    \begin{itemize}
        \item Se descartaron puntos con banderas: \texttt{pixelFlags\_nodata}, \texttt{bad}, \texttt{cr}, \texttt{edge}, \texttt{saturated} o \texttt{psfFlux\_flag}.
        \item Solo se conservaron curvas con $>5$ puntos válidos tras remover banderas.
        \item Se filtraron curvas con dispersión intra-época anómala (descartando así la banda $g$ en Fornax, la banda $y$ en 47Tuc y el campo Seagull).
    \end{itemize}
    \item \textbf{Ajuste de plantillas (Braga et al. 2024):}
    \begin{itemize}
        \item Ante la escasez de datos ($6$ a $\sim70$ puntos), los ajustes de Fourier libres divergen; se usó el código de plantillas \texttt{RRL\_lcvtemplate\_griz\_LSST}.
        \item El periodo $P$ es obligatorio. La época de fase máxima ($t_0$) se fijó en la banda $r$ y se ancló a los demás filtros.
        \item \textit{Aproximaciones:} Plantilla de $g$ para la banda $u$; plantilla de $z$ para la banda $y$. Para RRc no hay plantillas en banda $z$, por lo que no se ajustaron las bandas $z$ ni $y$ en RRc.
    \end{itemize}
    \item \textbf{Agrupamiento de fase (\textit{Phase Clumping}):} Cada época de DP1 consistió en 5 a 20 visitas sucesivas de 30 s tomadas en pocas horas. Al desfasar con el periodo ($P \sim 0.2 - 0.7$ d), las observaciones se concentran en racimos muy apretados, dejando valles o crestas desprovistos de datos.
\end{itemize}

\section{Sección 4: Campos Limpios y Poco Poblados (EDFS, Rubin95, Rubin38)}
Se analizaron en detalle las 4 estrellas con mejor cobertura fotométrica (Tabla 2 del artículo):
\begin{itemize}
    \item \textbf{Efecto Blazhko en EDFS:} Para la estrella \texttt{SSS J035520.3-484728} (RRab, $P = 0.5654$ d), la plantilla en banda $r$ no ajustaba bien la dispersión máxima. Al contrastar con la banda $G$ de Gaia DR3, se confirmó que es una RR Lyrae modulada por efecto Blazhko.
    \item \textbf{Metalicidad fotométrica $[\text{Fe/H}]_{gri}$ (Marconi et al. 2022):}
    \begin{equation*}
        [\text{Fe/H}]_{gri} = \alpha(g - r)_0 + \beta(r - i)_0 + \gamma \quad (\sigma = 0.31\text{ dex})
    \end{equation*}
    \begin{itemize}
        \item Para EDFS y las dos estrellas de Rubin95: $[\text{Fe/H}]_{gri} \sim -1.8$ dex, en excelente concordancia con trazas teóricas de baja metalicidad.
        \item \textbf{Alerta en Rubin38:} Un desfase de solo $0.14$ mag en la magnitud media de banda $i$ provocó valores espurios de $[\text{Fe/H}] = -3.59$ dex (\texttt{Object}) y $+1.41$ dex (\texttt{DiaObject}). Demuestra que la metalicidad fotométrica es \textbf{extremadamente sensible} al error fotométrico en cualquier banda.
    \end{itemize}
    \item \textbf{Módulos de Distancia ($\mu$) y el Índice de Wesenheit ($W_{gr}$):}
    \begin{equation*}
        W = r - 2.796(g - r)
    \end{equation*}
    $W$ es independiente del enrojecimiento por construcción y su término de metalicidad es despreciable ($\sim 0.05\text{ mag/dex}$). Los valores de $\mu_W$ coinciden perfectamente entre \texttt{Object} y \texttt{DiaObject} y con la literatura (ej. EDFS: $\mu_W = 17.22$ vs $\mu_{lit} = 17.19$).
    \item \textbf{Amplitudes en el Diagrama de Bailey (Figura 6):}
    \begin{itemize}
        \item Las amplitudes observadas coinciden con los \textbf{modelos ZAHB} (Rama Horizontal de Edad Cero).
        \item Son \textbf{totalmente inconsistentes con los modelos evolucionados} (\textit{evolved models}), los cuales predicen periodos desplazados y amplitudes erróneas.
    \end{itemize}
\end{itemize}

\section{Sección 5: Campos Abarrotados (47Tuc y Fornax)}
\begin{itemize}
    \item \textbf{Técnica de Razones Fijas de Amplitud:} Al disponer solo de 2 noches (Fornax) y 4 noches (47Tuc), se fijaron las amplitudes de $g$ e $i$ respecto a $r$ según las razones de Braga et al. (2024): $A_g/A_r$ y $A_i/A_r$.
    \item \textbf{47Tuc (SMC de fondo):} 34 RR Lyrae recuperadas en \texttt{diaObjectForcedSource} frente a solo 1 en \texttt{Object}. Nueve estrellas con datos $gri$ permitieron estimar $[\text{Fe/H}]_{gri}$, coincidiendo con la literatura de Li et al. (2023) cuando la cobertura de fase no fue deficiente.
    \item \textbf{Fornax (108 estrellas con fotometría simultánea):} La diferencia media entre \texttt{Object} y \texttt{DiaObject} fue de solo $\sim0.01$ mag en $r$ y $\sim0.03$ mag en $i$, demostrando la consistencia global del pipeline de Rubin.
    \item \textbf{Comparación de Módulos de Distancia:}
    \begin{itemize}
        \item La relación PWZ de Wesenheit ($W_{gr}$) reproduce la distancia de la SMC en 47Tuc con una desviación de apenas $\Delta\mu = 0.05$ mag.
        \item En cambio, las relaciones teóricas PLZ mono-banda ($\mu_r$ y $\mu_i$) arrojan distancias \textbf{sistemáticamente sobreestimadas en $+0.25$ a $+0.30$ mag} tanto en 47Tuc como en Fornax ($\Delta\mu_r = +0.28 \pm 0.22$, $\Delta\mu_i = +0.31 \pm 0.21$).
    \end{itemize}
\end{itemize}

\section{Sección 6: Discusión, Conclusiones y Recomendaciones}
\begin{enumerate}
    \item \textbf{Uso conjunto de productos fotométricos:} Aunque \texttt{diaObjectForcedSource} es prioritario en variables, los autores recomiendan fuertemente contrastarlo con \texttt{objectForcedSource} en series temporales tempranas como herramienta para detectar curvas defectuosas o ajustes divergentes.
    \item \textbf{Superioridad del Índice de Wesenheit ($W_{gr}$):}
    Al comparar todos los módulos de distancia $\mu_W$ con la literatura (Figura 11 del paper), el acuerdo es sobresaliente:
    \begin{equation*}
        \Delta\mu_W = \mu_W(\text{DP1}) - \mu_{lit} = -0.01 \pm 0.36\text{ mag}
    \end{equation*}
    Se recomienda formalmente emplear la relación PWZ en $W_{gr}$ para futuros trabajos de escala de distancias en LSST.
    \item \textbf{Causa del sesgo sistemático en relaciones PLZ mono-banda:}
    El exceso de $+0.25$ a $+0.30$ mag en las relaciones PLZ se debe a que Marconi et al. (2022) incluyeron \textit{modelos evolucionados} al calibrar sus ajustes teóricos. Los modelos evolucionados tienen luminosidades intrínsecas $0.2\text{ dex}$ más brillantes que las ZAHB, inflando artificialmente el módulo de distancia ($\mu = m - M$).
    \item \textbf{Perspectiva futura:} Es prioritario que los teóricos recalibren las relaciones PLZ para LSST excluyendo modelos evolucionados incompatibles con las amplitudes observadas. Se sugiere calibrar empíricamente con cúmulos globulares monometálicos observados por Rubin.
\end{enumerate}

\end{document}
